\begin{frame}{실습 노트북 지도 + 파이프라인 요약}
  \begin{itemize}
    \item \textbf{§1} 스팸 데이터 만들고 스케일링
    \item \textbf{§2} 본문 \texttt{logistic\_gradient\_descent}를
      numpy로 학습 (그래디언트 $(h-y)\cdot \boldsymbol{x}$ 확인)
    \item \textbf{§3} $\sigma'(z) = \sigma(z)(1-\sigma(z))$ 를
      수치로 검증
    \item \textbf{§4} 임계값 sweeping $\to$ P/R 트레이드오프와 비용
      최소화 임계값 ($\approx 0.14$) 직접 계산
    \item \textbf{§5} PR-AUC/ROC-AUC와 ``무조건 0'' 모델의 기준선
      (= 양성 비율) 검증
    \item \textbf{§6} 본문 실습 코드 실행
  \end{itemize}
  \vskip0.6em
  \begin{block}{파이프라인 요약}
    선형회귀와 로지스틱회귀는 겉보기엔 다른 문제(숫자 예측 vs
    확률 예측)를 풀지만, 뜯어보면
    \kb{``선형모델 $\to$ 손실함수 정의 $\to$ 경사하강법으로
    최소화''}라는 \textbf{완전히 같은 파이프라인} ---
    이 패턴은 \textbf{이번 학기 내내 반복}된다.
  \end{block}
\end{frame}
